\documentclass[11pt]{article} %inicia el documento con una determinada clase
\usepackage{graphicx} %permite incluir imágenes
\usepackage[utf8]{inputenc} %permite utilizar caracteres extras
\usepackage[spanish]{babel} %modifica en lenguage del documento
\usepackage[vmargin=4cm,tmargin=3cm,hmargin=2cm,letterpaper]{geometry} %modifica lor margenes
\setlength{\parindent}{0pt} %permite modificar la indentanción
\usepackage{fancyhdr}
\pagestyle{fancy}
\renewcommand{\headrulewidth}{0.5mm} %le da un grozor a la linea del encabezado
\pagenumbering{gobble} %quita la numeración de página
\graphicspath{{images/}}
\usepackage{framed,color}
\definecolor{shadecolor}{rgb}{1,0.8,0.3}

\begin{document}
\fancyhead[L]{\includegraphics[scale=0.15]{ucr}}
\fancyhead[C]{UNIVERSIDAD DE COSTA RICA\\ESCUELA DE INGENIERÍA ELECTRICA\\PROGRAMACIÓN BAJO PLATAFORMAS ABIERTAS\\PROYECTO DE DESARROLLO - BOLETIN.02}
\fancyhead[R]{\includegraphics[scale=0.4]{eie}}
\pagestyle{fancy}
\setlength{\headheight}{30pt}
\rule{\linewidth}{0.2mm}
\vspace{10pt}

\paragraph{Primer comunicado:}
Por favor usen el git de la escuela, lo que podemos hacer es tener
las cosas pesadas (que está en /extra/pdf) lo usemos en gitHUb y la del programa en si y eso en el git de la escuela. Esto por que hablando con el profe el quiere que usemos el de la escuela. Le comunique que teníamos problemas y di me dijo "eso es parte de lo que quiero que experimenten".. entonces.. a arreglar el git, A mi ya me funciona entonces ahí traten de arreglarlo, si no ahí vemos.

\paragraph{Segundo comunicado:}
Ya masomenos sabemos que hace cada uno. Lo pondré en la siguiente nota
\begin{itemize}
\item Luis Diego: va a ordenar la tabla por medio de regular expressions, esto para automatizar la cosa. Como creo que lo va hacer en python o pear, que suba al git el scrip documentado para que todo entendamos.\\
\begin{shaded}
\begin{center}
	TABLA 1: Estudiantes\\
    \begin{tabular}{ | l | l |}
    \hline
    id & carné \\ \hline
    AUTO-INCREMENT & XXXXXX   \\
    \hline
    \end{tabular}
\end{center}

\begin{center}
	TABLA 2: Materias\\
    \begin{tabular}{ | l | l |}
    \hline
    id & sigla \\ \hline
    AUTO-INCREMENT & XXXXXX \\ 
    \hline
    \end{tabular}
\end{center}

\begin{center}
	TABLA 3: GENERAL\\
    \begin{tabular}{ | l | l || l | l | l | l | l |}
    \hline
    carné & recinto & año & periodo & grupo & sigla & nota \\ \hline
    LINK-TO-1ST-TABLE & XX & XXXX & X & XXX & LINK-TO-2ND-TABLE  & X.X \\
    \hline
    \end{tabular}
\end{center}
\textit{nota: vean que las dos primera tablas tienen un id (eso es el primary key, si no lo han escuchado en cualquier libro de bases de datos debe salir) y con ese número interconectan las tablas)}
\end{shaded}
\item Sebastian: se dedicará a hacer los queries a las tablas por medio de php (ya masomenos viendo las tablas de arriba va viendo como van a estan configuradas las tablas (en el boletin sale como hacer las tablas (o bajese phpmyadmin, es más fácil) , puede hacer 3 de prueba que cumplan con las tablas de arriba y mandarse a hacer los filtros
\item Carla: no se si vos querés darle con php o bien dedicarte a lo de hacer el reporte, (no se como estará sebas con lo de él) pero almenos lo del reporte di hay q ver como se hacen los histográmas como se pasan a un .tex y así, que a eo me pienso a dedicar. Aun no se ni papa de hacer gráficos y eso. =P
\end{itemize}
\end{document}